%% derivation.tex
\section{From Mole Balances to the Flux Constraint}\label{sec:derivation}

Flux balance analysis (FBA) uses an algebraic constraint derived from a
mole balance. Writing the balance out identifies the assumptions used
to obtain its familiar textbook form and the operating conditions under
which that reduction applies. Consider a well-mixed bioreactor of volume $V$ holding a growing cell
culture. The culture exchanges mass with its surroundings through a set of
physical streams $s$ (feed, harvest, gas sparging) and hosts an
intracellular reaction network with $n$ fluxes $\hat v_j$, $j=1,\dots,n$,
acting on $m$ metabolite species. For a species $i$ present at
concentration $C_i$, an open mole balance over the whole culture accounts
for material entering and leaving through the streams together with
material produced or consumed by the reaction network:
\begin{equation}\label{eq:mole-balance}
  \sum_{s} d_s\, C_{i,s}\,\dot V_s \;+\; \sum_j \sigma_{ij}\hat v_j\, V
  \;=\; \frac{d}{dt}\bigl(C_i V\bigr),
\end{equation}
where $d_s=+1$ for streams entering the reactor and $d_s=-1$ for streams
leaving it, $C_{i,s}$ is the concentration of species $i$ in stream $s$,
$\dot V_s$ is the volumetric flow rate of that stream, and $\sigma_{ij}$ is
the stoichiometric coefficient of species $i$ in reaction $j$, the $(i,j)$
entry of the stoichiometric matrix $\mathbf S\in\mathbb R^{m\times n}$.
This mole balance, Equation~\eqref{eq:mole-balance}, does not require
assumptions about growth rate, feeding strategy, or steady state and
provides the starting point for the reduction below.

Two steps turn Equation~\eqref{eq:mole-balance} into a statement about
fluxes alone. First, write the reactor volume in specific units,
$V=B\bar V$, with $B$ the total biomass in the reactor and $\bar V$ the
culture volume per unit biomass; this ties the extensive balance to the
biomass that is actually growing, and the specific growth rate follows
as $\mu\equiv(1/B)\,dB/dt$. Second, the product rule expands the
accumulation term as:
\begin{equation}\label{eq:product-rule}
  \frac{d}{dt}\bigl(C_iV\bigr)=\frac{d}{dt}\bigl(C_iB\bar V\bigr)
  = B\bar V\,\frac{dC_i}{dt}
  + C_i\bar V\,\frac{dB}{dt}
  + C_iB\,\frac{d\bar V}{dt}.
\end{equation}
Dividing Equation~\eqref{eq:mole-balance} through by $V=B\bar V$ and
substituting Equation~\eqref{eq:product-rule} isolates the intensive
balance for species $i$:
\begin{equation}\label{eq:intensive-balance}
  \frac{1}{B\bar V}\sum_{s} d_s\, C_{i,s}\,\dot V_s
  \;+\; \sum_j \sigma_{ij}\hat v_j
  \;=\; \frac{dC_i}{dt} \;+\; \mu\,C_i
  \;+\; \frac{C_i}{\bar V}\,\frac{d\bar V}{dt},
\end{equation}
where the middle term on the right, $\mu C_i$ with
$\mu=(1/B)\,dB/dt$, is the growth-dilution term produced by the product
rule. Three assumptions each remove one of the remaining terms. If the
intracellular pool of species $i$ is at pseudo-steady state
($dC_i/dt=0$, the usual separation-of-timescales argument, since
metabolite turnover on the order of seconds to minutes is fast relative
to a doubling time of hours), if the specific culture volume is constant
($d\bar V/dt=0$, so the working volume tracks the growing biomass rather
than drifting relative to it), and if species $i$ crosses no physical
stream (the sum over $s$ vanishes for this species), the intensive
balance reduces to the growth-dilution constraint:
\begin{equation}\label{eq:dilution-constraint}
  C_i\,\mu \;=\; \sum_j \sigma_{ij}\hat v_j
  \qquad\Longleftrightarrow\qquad
  \boxed{\mathbf S\hat{\mathbf v} \;=\; \mu\mathbf x},
\end{equation}
where $\mathbf x=(C_1,\dots,C_m)^\top$ collects the intracellular
concentrations and $\hat{\mathbf v}\in\mathbb R^n$ collects the fluxes.
The resulting flux constraint, Equation~\eqref{eq:dilution-constraint},
retains growth dilution: net metabolic production of each intracellular
species equals its dilution rate in the growing culture.

The form used throughout the constraint-based modeling literature is
obtained by neglecting the right-hand side of
Equation~\eqref{eq:dilution-constraint}:
\begin{equation}\label{eq:sv0}
  \mathbf S\hat{\mathbf v} \;=\; \mathbf 0.
\end{equation}
The steady-state relation in Equation~\eqref{eq:sv0} is a special case of
Equation~\eqref{eq:dilution-constraint}, rather than an independent
starting point \cite{Orth2010}. It is an appropriate approximation when the two
conditions that justify neglecting $\mu\mathbf x$ hold together: growth is
slow relative to metabolic turnover, and intracellular pools are dilute
enough that $\mu\mathbf x$ stays small in absolute terms even when $\mu$
itself is not negligible. Both conditions can fail in relevant operating
regimes. In fed-batch culture the feed is metered to increase
the working volume over the course of a run, so the specific volume
drifts ($d\bar V/dt\neq0$); the assumption used to reach
Equation~\eqref{eq:dilution-constraint} no longer holds, and the omitted
term carries information about the feed schedule that
Equation~\eqref{eq:sv0} cannot see. In cell-free systems there is no
growing cell to dilute into ($\mu=0$), yet the reaction mixture remains
time varying: transcript, protein, and metabolite pools
accumulate and decay over the course of a batch reaction ($d\mathbf
x/dt\neq0$), so the pseudo-steady-state assumption does not apply. A
cell-free capstone later in this chapter returns to this setting, where
the accumulation term omitted by Equation~\eqref{eq:sv0} is measured
directly.

The dilution-retaining constraint and its steady-state special case,
Equations~\eqref{eq:dilution-constraint} and~\eqref{eq:sv0}, are, as
written, statements about a closed system: every species that appears has its production and consumption
accounted for entirely by the internal network $\mathbf S\hat{\mathbf v}$,
with the physical stream terms of Equation~\eqref{eq:mole-balance}
already omitted along the way. Cellular metabolism is open, since cells
take up substrates and secrete products across the boundary of whatever
system is being modeled. Constraint-based reconstructions restore that
openness without reintroducing the stream terms explicitly: hypothetical
exchange, or boundary, reactions are added to the network itself, one
pseudo-reaction $\emptyset\rightleftharpoons M_i$ for each boundary
metabolite $M_i$, entering $\mathbf S$ as an isolated nonzero column with
no reaction partner \cite{Bordbar2014}. Flux through an exchange reaction
enters Equation~\eqref{eq:dilution-constraint} or
Equation~\eqref{eq:sv0} exactly as any metabolic flux does, but it is now
interpretable as uptake or secretion across the system boundary rather
than as a step of intracellular chemistry. A closed algebraic statement
can thereby represent a genuinely open system: openness becomes entirely
a matter of which reactions are admitted into $\mathbf S$ and how their
bounds are set, the question taken up next.
